\documentclass[10pt]{article}
\usepackage[margin=20mm]{geometry}
\usepackage{amsmath,amssymb,mathtools,fontspec,unicode-math,xurl,hyperref}
\IfFontExistsTF{DejaVu Serif}{\setmainfont{DejaVu Serif}}{\setmainfont{Latin Modern Roman}}
\setmathfont{latinmodern-math.otf}
\setlength{\parindent}{0pt}\setlength{\parskip}{6pt}\sloppy
\title{Formal initial review: M25-0073
Minimal hypersurfaces for generic metrics in dimension 8}
\author{}
\date{26 October 2026}
\begin{document}
\maketitle
\section*{Scope and outcome}
The primary object is the acquired 111-page publisher PDF. DOI:\par
https://doi.org/10.1007/s00222-025-01333-0.\par
All 111 pages were read as extracted text. Thirty-one selected source pages were additionally inspected as actual page images. The associated source tree is the older author version arXiv:2205.01047v1 (2 May 2022); it is not the exact source of the journal PDF. All four findings below are established directly against the publisher PDF.\par
Initial outcome: Major. The global compactness assertion used to obtain the countable covering in Theorem 3.4 lacks uniform positivity of the ambient metrics. A sequence satisfying the stated bounds and remaining outside the excluded higher-multiplicity neighborhood disproves that assertion. This is a demonstrated failure of a step in the principal proof, even though a countable ellipticity exhaustion repairs this step. It does not disprove Theorem 1.2. Three further statements have local scope defects for which no damage to the main conclusion was established.\par
No. | Confirmed issue in the bound version | Severity / scope\par
1 | The global bounded family need not have a Riemannian limit | Major: countable-cover argument\par
2 | A two-sided lift does not determine a unique ambient double cover | Minor: preliminary assertion\par
3 | The first Dirichlet eigenspace need not be simple on disconnected \(\mathcal U\) | Minor: Lemma D.3(i)\par
4 | Lemma D.6 includes nonsingular invertible examples for which its set is empty | Minor: auxiliary lemma scope\par
The main theorem concerns connected embedded locally stable minimal hypersurfaces (LSMHs), with its specified area and singular-set conditions. It is not read here as a claim about arbitrary stationary varifolds. The paper’s one-sided index convention is the index of the two-sided cover. These conventions were retained.\par
Page convention. Locations below give printed journal pages; the one-based PDF page is the printed number minus 1192. Verification describes only the stated check, not all arguments of the paper. No issue quota was imposed.\par
\section*{Problem 1. The exhaustion used for global compactness omits uniform positivity}
Location and exact claim. Section 9.1, p. 1253 (PDF 61), defines \[\mathcal M^{k,\alpha}(M;\Lambda,I)=
 \{(g,\Sigma)\in\mathcal M^{k,\alpha}(M):
 \|g\|_{C^{k,\alpha}}\leq\Lambda,\ 
 \mathcal H^7_g(\Sigma)\leq\Lambda,\ \operatorname{index}(\Sigma,g)\leq I\}.\] The set \(\mathcal M_{>1}^{k,\alpha}(M;\Lambda,I,\delta)\) consists of pairs in these bounds within \(\delta\), in the metric-factor \(C^{k-1,\alpha}\) distance plus varifold \(\mathbf F\) distance, of \((\widetilde g,m|\widetilde\Sigma|)\) for some bounded pair and integer \(m\geq2\). The text asserts that its complement in the bounded family is compact for every \(\delta>0\), citing Theorem G.1, and uses this for the finite cover (76). Theorem G.1, p. 1299 (PDF 107), assumes \(g_i\to g\) in \(C^3\) with \(g\) a Riemannian metric.\par
Counterexample: all geometric assumptions. Let \(M=S^1_{2\pi}\times T^7\), where \(T^7=\mathbb R^7/\mathbb Z^7\) has its flat unit-volume metric \(h\). Set \[\rho(s)=2-\cos s,\qquad
 g_* = ds^2+\rho(s)^2h,\qquad \Sigma=\{0\}\times T^7 .\] This is a smooth closed eight-manifold and a smooth connected embedded two-sided hypersurface. Since \[\rho(0)=1,\quad \rho'(0)=0,\quad \rho''(0)=1,
 \qquad A_\Sigma=0,\qquad
 \operatorname{Ric}_{g_*}(\partial_s,\partial_s)|_\Sigma=-7,\] \(\Sigma\) is minimal, has area one, and its Jacobi operator is \(L=\Delta_h-7\). In the article’s convention, \[Q(u,u)=\int_{T^7}(|\nabla u|^2+7u^2)\,d\mu_h>0
 \quad (u\ne0).\] It is strictly stable, hence locally stable and of index zero. Its singular set is empty; all stated measure conditions are satisfied.\par
Put \(H(t)=e^{-1/t}\) for \(t>0\), \(H(t)=0\) for \(t\leq0\), and \[\chi(s)=\frac{H(\cos s-\tfrac12)}
 {H(\cos s-\tfrac12)+H(\tfrac34-\cos s)},\quad
 q_\varepsilon=\chi+\varepsilon(1-\chi),\quad
 g_\varepsilon=q_\varepsilon g_*,\quad 0<\varepsilon\leq1.\] The denominator is positive everywhere. The smooth function \(\chi\) equals one on a fixed neighborhood of \(0\) and zero near \(\pi\). Every \(g_\varepsilon\) is positive definite and equals \(g_*\) near \(\Sigma\). Thus \(\Sigma\) retains area one and index zero.\par
For any fixed background tensor norm the \(C^{k,\alpha}\) norms of \(g_\varepsilon\) are uniformly bounded, because \(g_\varepsilon=\chi g_*+\varepsilon(1-\chi)g_*\). Choose an integer \(\Lambda\) above this bound and above one, and \(I=1\). All pairs satisfy the displayed bounds. Yet \(g_{1/j}\to\chi g_*\) in \(C^\infty\), and the limit vanishes on an open set. No subsequence has a Riemannian limit, even in \(C^0\).\par
\section*{Problem 1 continued: the excluded set and the repair}
The higher-multiplicity exclusion does not remove the sequence. There is a single \(\delta_0>0\) for which all pairs \((g_\varepsilon,\Sigma)\), \(0<\varepsilon\leq1\), avoid \(\mathcal M_{>1}^{k,\alpha}(M;\Lambda,I,\delta_0)\). Otherwise choose \(\varepsilon_j\), metrics \(\widetilde g_j\), connected LSMHs \(\widetilde\Sigma_j\), and integers \(m_j\geq2\) with \[\|\widetilde g_j-g_{\varepsilon_j}\|_{C^{k-1,\alpha}}
       +\mathbf F(m_j|\widetilde\Sigma_j|,|\Sigma|)\longrightarrow0.\] On a fixed tubular neighborhood \(U\) of \(\Sigma\) all \(g_{\varepsilon_j}=g_*\). Hence \(\widetilde g_j\to g_*\) in \(C^3\) on \(U\), while the integral stationary varifolds \(V_j=m_j|\widetilde\Sigma_j|\) converge there to the smooth multiplicity-one varifold \(|\Sigma|\).\par
Choose a regular point \(x\) of \(\Sigma\). The local geometric bounds are uniform. Local monotonicity gives upper semicontinuity of densities at points \(x_j\to x\). Such \(x_j\) can be chosen in the regular part of \(\widetilde\Sigma_j\), since the limit has positive mass in every ball about \(x\). Consequently, \[2\leq m_j=\Theta(V_j,x_j),\qquad
 \limsup_j\Theta(V_j,x_j)\leq\Theta(|\Sigma|,x)=1,\] a contradiction. Equivalently, Allard regularity at the smooth multiplicity-one limit excludes such integer multiplicities. Only the fixed positive geometry on \(U\) is used; degeneration outside \(U\) is irrelevant.\par
Take \(\delta=\delta_0/2\) (or \(1/J<\delta_0\)). The entire noncompact sequence belongs to the precise complement claimed to be compact. This verifies the exclusion condition rather than merely exhibiting arbitrary degenerating metrics.\par
Counterevidence. The local canonical neighborhoods in Definition 3.1, pp. 1211–1212, also impose closeness to a fixed positive metric. That condition supplies positivity and does not suffer from this example. Remark 3.1 uses the coarser \(C^{k-1,\alpha}\) topology; the example fails even there. No convention incorporating a bound on \(g^{-1}\) was found in the global definition or surrounding argument. Such a convention or an explicit bound is needed.\par
Impact. The invalid inference precedes (76), the countable cover (77), the refinement (78), and Theorem 3.4 on p. 1255. Theorem 3.4 and Theorem 3.3 yield Corollary 3.5 (pp. 1213–1214), used in Theorem 1.2 (pp. 1215–1219). This explicit dependency supports Major under the prescribed calibration. It is a defect in the supplied derivation, not a refutation of generic regularity.\par
Repair and limit. Fix a smooth reference metric \(g_{\rm ref}\) and add \(g\geq L^{-1}g_{\rm ref}\) to each bounded family, \(L\in\mathbb N\). The metric limit is now positive. Bounded \(C^{k,\alpha}\) norms give \(C^{k-1,\alpha}\) precompactness, and Theorem G.1 applies. A higher-multiplicity limit would put the original pairs in the excluded \(\delta\)-set eventually, so the complementary stratum is compact. Every positive metric on compact \(M\) satisfies such a bound for some integer \(L\); this extra countable index preserves countability. This repairs the identified compactness step. It is not a reproof of all later tree refinements or the main theorem.\par
\section*{Problem 2. The ambient double cover is not unique}
Location and claim. Section 2, p. 1200 (PDF 8), immediately before (1), asserts that for a connected one-sided \(\Sigma\subset M\) there is a unique ambient double cover \(\widehat M\to M\) whose inverse image of \(\Sigma\) is connected and two-sided.\par
Complete example. Take \[M=\mathbb{RP}^8\#(S^1\times S^7),\qquad
 \Sigma=\mathbb{RP}^7\subset\mathbb{RP}^8,\] with the connected sum performed in a ball disjoint from \(\Sigma\). Choose a metric round near this standard projective hyperplane and extend it smoothly across the connected sum. The hypersurface is closed, connected, smooth, embedded, one-sided and totally geodesic. It is minimal and locally stable: on sufficiently small intrinsic balls the Dirichlet gradient term dominates the bounded second-variation potential. At points off \(\Sigma\), small balls miss it. Its area is finite and its singular set empty, so it is an LSMH.\par
Van Kampen gives \(\pi_1(M)=\langle a,b\mid a^2=1\rangle=\mathbb Z_2*\mathbb Z\). The two characters \[\chi_0(a)=1,\quad\chi_0(b)=0,\qquad
 \chi_1(a)=1,\quad\chi_1(b)=1
 \quad\text{in }\mathbb Z_2\] are surjective and have distinct kernels. They define inequivalent connected double covers over \(M\). Both restrict nontrivially to \(\pi_1(\Sigma)=\langle a\rangle=\mathbb Z_2\). For either cover the full inverse image of \(\Sigma\) is its connected cover \(S^7\). The normal line bundle of \(\mathbb{RP}^7\) in \(\mathbb{RP}^8\) has nontrivial orientation character on \(a\), which is killed on \(S^7\); hence both lifts are two-sided. Pullback metrics preserve minimality and local stability. The lift of a \(b\)-loop distinguishes the ambient covers.\par
Counterevidence, impact and repair. The intrinsic normal orientation double cover of \(\Sigma\) is canonical. That does not make the ambient cover unique. Subsequent sign and lifted Jacobi arguments need an appropriate cover, not uniqueness among all ambient covers. No failure of those arguments was established. Specify a choice, for example the cover determined by the mod-two class dual to \(\Sigma\), or state the canonical property for the intrinsic cover alone. The example does not dispute existence. The impact is local; severity Minor.\par
\section*{Problem 3. Simplicity of the first eigenspace needs connectedness}
Location. Appendix D allows an arbitrary precompact open \(U\subset M\) with \(\partial U\cap\operatorname{Sing}(\Sigma)=\varnothing\) (p. 1273, PDF 81), and sets \(\mathcal U=U\cap\Sigma\). Lemma D.3(i), p. 1275 (PDF 83), asserts \(\dim E_1=1\) and a positive first eigenfunction, without requiring \(\mathcal U\) connected.\par
Use \((M,g_*,\Sigma)\) from Problem 1. Choose \(U\) as two disjoint small product neighborhoods whose intersections with \(T^7\) are isometric Euclidean balls \(B_r^{(1)},B_r^{(2)}\), and take the lemma’s potential \(h=0\). All stated conditions hold. The Dirichlet space and operator split as a direct sum. If \(\phi\) is the first Dirichlet eigenfunction of \(-\Delta+7\) on \(B_r\), then \((\phi,0)\) and \((0,\phi)\) are independent first eigenfunctions with the same lowest eigenvalue. Thus \(\dim E_1=2\).\par
Require \(\mathcal U\) connected for simplicity, or state the result componentwise. The main uses on connected \(\Sigma\) and near a single isolated conical singularity do not yield a demonstrated central failure from this example. The spectral decomposition itself survives. Severity Minor.\par
\section*{Problem 4. Lemma D.6 must exclude the nonsingular invertible case}
Location and quantifiers. Lemma D.6, p. 1283 (PDF 91), takes a two-sided LSMH and a finite group \(G\) acting freely on \(M\), preserving the geometric data. It claims that \[\mathcal F_G^{k,\alpha}=
 \{f\in C^{k,\alpha}(M)^G:
 L_{\Sigma,g}u=\nu(f)
 \text{ has no slower-growth $G$-invariant solution}\}\] contains an open dense subset of \(C^{k,\alpha}(M)^G\). There is no assumption \(\operatorname{Sing}(\Sigma)\ne\varnothing\).\par
Counterexample and all hypotheses. Use the same smooth closed warped product \((M,g_*,\Sigma)\) from Problem 1 and \(G=\{\mathrm{id}\}\). This is a finite free action, exactly the trivial-group case also used on p. 1291. The hypersurface is two-sided and an LSMH. Its Jacobi operator is \(L=\Delta_h-7\) and is invertible. For every \(f\in C^{k,\alpha}(M)\) set \(b=\nu(f)|_\Sigma\). In Fourier variables on \(T^7\) the unique solution is \[\widehat u(m)=-\frac{\widehat b(m)}
                      {4\pi^2|m|^2+7},\qquad m\in\mathbb Z^7.\] Elliptic regularity gives the required local regularity and \(u\in L^2(\Sigma)\). Since \(\operatorname{Sing}(\Sigma)=\varnothing\), Definition 2.5 on p. 1206 (PDF 14) explicitly says that every \(L^2\) function satisfies its slower-growth condition. Therefore every \(f\) has such a solution: \(\mathcal F_{\{\mathrm{id}\}}^{k,\alpha}=\varnothing\). It cannot contain an open dense subset of the nonempty function space.\par
Counterevidence and core use. The proof of D.6 constructs functions at singular points. Its intended application is Theorem 2.6, which explicitly assumes a nonempty singular set on p. 1206. That theorem’s proof on p. 1291 (PDF 99) invokes D.6 under the same assumptions. Thus the example invalidates the wider lemma but does not invalidate that application. Add \(\operatorname{Sing}(\Sigma)\ne\varnothing\) to D.6. This check validates the necessity and compatibility of the restriction; it does not reconstruct every step of the restricted lemma. Severity Minor.\par
\section*{Candidates not counted as scientific findings}
The actual image of Lemma 7.1, p. 1228 (PDF 36), shows condition (b) ending after the word and without the expected density display. The older v1 source, lines 1370–1374, supplies density pinching and equality of the two singular densities. That older formula is version-specific explanatory evidence, not silently inserted into the journal text. Density assumptions are also visible in later journal statements. This typesetting omission is not counted under the assigned scientific scope.\par
There is a squared-norm convention mismatch between the integral in (32), p. 1222, and the norm scaling in (44), p. 1229. Taking the square root restores the scaling, and monotonicity is unchanged under the positive square root. Author v1 explicitly defines the square of \(J\) as the integral in its local growth lemma. This notational matter is not promoted to a further scientific defect. Unresolved questions about singular sheeting constants or external min–max applications remain coverage limits.\par
\section*{Independent positive checks and proof-chain audit}
Logarithmic cone mode (Appendix A). For the repeated-root case, write \(a=2\alpha'\), \(X=K^a\), \(L_0=\log K\), and \(z=\log r\). Evaluating the defining integral gives a positive factor \(r^a\) times \[Az^2+2Bz+C,\qquad
 A=\frac{(X-1)^3}{a},\quad
 B=\frac{3XL_0(X-1)^2}{a}-\frac{(X-1)^3}{a^2}.\] The independently expanded finite difference and determinant give \[B^2-AC=\frac{(X-1)^4}{a^2}
 \left(3XL_0^2-\frac{(X-1)^2}{a^2}\right).\] Both identities were checked symbolically against the defining integral, without running author code. With \(x=a\log K\), the sign condition is \(3x^2<4\sinh^2(x/2)\). It holds for \(|x|\geq6\): \(\sinh t/t\) increases for \(t>0\), and \(\sinh3>3+3^3/6>3\sqrt3\). Thus \(K\geq\max(3,e^{3/\sigma})\) suffices when \(|a|\geq2\sigma\). This verifies the uniform large-\(K\) step for the logarithmic mode, including negative \(a\), not every nonlinear growth estimate.\par
Critical compactness barrier (Appendix E). At \(n=7\), the two powers in the radial barrier are \(-3\) and \(-49/24\), exactly the roots of \[a(a-1)+\frac{145}{24}a+\frac{49}{8}=0.\] The error margin \(1/24-1/196-1/70=131/5880\) is positive. The independent script verifies these rational identities. The critical Hardy case was not dismissed by assuming all tangent cones strictly stable.\par
Lipschitz Sard argument. Lemma 8.5, pp. 1243–1244 (PDF 51–52), defines critical points by all limiting normalized secants landing in a proper subspace. Lipschitz boundedness makes the image of a sufficiently small relative ball lie in an arbitrarily thin slab around that subspace. Its \(n\)-volume is at most a constant times slab thickness times \(r^n\). The disjoint-ball covering and fivefold enlargement in the proof yield zero image measure as thickness tends to zero. This supports the measure step; it does not alone prove the geometric slice parametrization.\par
Main dependency and boundary cases. The definitions of slower growth, semi-nondegeneracy, canonical neighborhoods and the one-sided index were traced into Sections 3–4. The index induction uses the integer-valued singular capacity, finitely many maximizers, and a strict capacity drop. The four multiplicity/index alternatives on p. 1218 were read individually: a higher-multiplicity limit has index zero by G.1; regular strong nondegeneracy excludes certain limits; the singular stable case uses the strict capacity drop. The countable intersection over area and index bounds on p. 1219 is compatible with these quantifiers, conditional on the component theorems.\par
The positive-Jacobi-field argument on p. 1300 was checked using \(Q(gh,gh)=\int h^2|\nabla g|^2\) for a positive solution \(h\). The one-sided case uses the cover and the paper’s index convention. Problem 1 is an applicability defect for G.1, not a counterexample to G.1 under its own assumptions.\par
\section*{Supporting checks}
Material | Actual coverage and limit\par
Sections 1–2 | Full text; scope, LSMH, index, cone exponents, slower growth; direct regular boundary-case checks.\par
Sections 3–4 | Full text; canonical topology, covering dependency, conformal perturbations, index/capacity induction.\par
Section 5 | Full text; application hypotheses/dependencies. External min–max results not independently reproved.\par
Sections 6–8 | Full text; growth normalizations, translated-cone rate, induced fields, slices and Lipschitz Sard. Long blow-up arguments not all reconstructed.\par
Section 9 | Full text; tree decomposition/covering chain. Complete global-compactness counterexample and local repair.\par
Appendices A–B | Full text; repeated-root algebra; graph formula dependence at statement level.\par
Appendices C–G | Full text; capacity recursion, spectral domains, critical barrier, deformation and compactness hypotheses. Singular sheeting not fully reconstructed.\par
Author source | Complete static parse of v1 tree; direct reading of lines 1351–1460 and 2221–2340. No author code or author TeX executed.\par
Evidence identity. The primary PDF has 111 pages, 2,340,175 bytes, SHA-256:\par
`c6f917e1f96406d0d820e38318efe5f1abcbe487a93fc2f5388085d961dee057`.\par
The older v1 `Ver1.0.tex` has 284,225 bytes, SHA-256:\par
`20a3d33d28b9aaf9b622dc141c179ad0e860993b3dc99d1110c9a70eae4b5a6a`.\par
All acquired files were checked against the manifest; full identities are in provenance.\par
Publisher metadata and arXiv history were inspected. Targeted publisher/arXiv correction searches did not establish an applicable correction. This limited negative result does not prove absence of corrections. Stable sources are the DOI and https://arxiv.org/abs/2205.01047.\par
Actual computation and compilation. Those checks and their unchanged evidence are retained. The cache snapshot is a conservative superset, not an exact consumed-file trace; previous build and page records remain in revisions.\par
Conclusion item | State\par
Initial judgment | Major: one core proof-step defect; three local statement defects.\par
Repair status | Ellipticity repairs the identified step; local scope corrections described. Full paper not recertified.\par
Author adoption | Not asserted; authors not contacted.\par
Remaining science | Complete singular sheeting/deformation reconstruction, external min–max inputs, downstream implementation of compactness repair.\par
Acceptance state | `human\_reviews=[]`; `scientific\_decision=pending`.\par
The current recovery performs the new compilation, report-page inspection and binding work. Administrative assistance had no source access or scientific role.\par
\end{document}
